% Diagrama: Pipeline de producción cartográfica de Transport-gis
% Usado en: 4-Capas-Codigo/4-3-Visualizacion.tex  (fig:tg-pipeline)
\begin{tikzpicture}[
  srcbox/.style = {rectangle, rounded corners=3pt, draw=gray!55, fill=gray!8,
                   minimum width=3.0cm, minimum height=0.90cm,
                   align=center, font=\small},
  extbox/.style = {rectangle, rounded corners=3pt, draw=unamAzul, fill=unamAzul!10,
                   minimum width=3.2cm, minimum height=0.90cm,
                   align=center, font=\small},
  intbox/.style = {rectangle, rounded corners=3pt, draw=unamAzul!60, fill=unamGrisClaro,
                   minimum width=6.4cm, minimum height=1.05cm,
                   align=center, font=\small},
  outbox/.style = {rectangle, rounded corners=3pt, draw=unamOro, fill=unamOro!15,
                   minimum width=3.2cm, minimum height=0.90cm,
                   align=center, font=\small},
  ->, >=Stealth, thick
]
%% — Fuentes externas —
\node[srcbox] (gtfs)  at (0.0, 9.2) {\gtfs{} SEMOVI~2024};
\node[srcbox] (inegi) at (5.4, 9.2) {INEGI~2023 (cartografía)};

%% — Componentes externos —
\node[extbox] (api) at (0.0, 7.5) {\textbf{Apimetro}\\{\footnotesize API REST (Go)}};
\node[extbox] (vft) at (5.4, 7.5) {\textbf{VFTModel}\\{\footnotesize (Python / NetworkX)}};

%% — Transport-gis: análisis —
\node[intbox] (scripts) at (2.7, 5.5)
  {{\footnotesize\texttt{apimetro\_client.py} $\cdot$ \texttt{vft\_runner.py}}\\
   {\footnotesize\texttt{pipeline/run\_pipeline.py}}};

%% — Transport-gis: datos —
\node[intbox] (gpkg) at (2.7, 4.0)
  {{\small\texttt{data/processed/}}\\
   {\footnotesize 10 capas \texttt{.gpkg} (Git~LFS)}};

%% — Transport-gis: QGIS —
\node[intbox] (qgis) at (2.7, 2.5)
  {{\small QGIS~3.x — \texttt{maps/projects/*.qgz}}\\
   {\footnotesize\texttt{maps/templates/*.qpt} (paletas UNAM)}};

%% — Transport-gis: exports —
\node[intbox] (exports) at (2.7, 1.0)
  {{\small\texttt{maps/exports/}}\\
   {\footnotesize\texttt{*.pdf} $\cdot$ \texttt{*.png}}};

%% — Destino final —
\node[outbox] (tesis) at (2.7, -0.5)
  {\textbf{Tesis LaTeX}\\{\footnotesize\texttt{Figures/}}};

%% — Flechas: fuentes → externos —
\draw[color=gray!60] (gtfs)  -- (api);
\draw[color=gray!60] (inegi) -- (vft);

%% — Flechas: externos → scripts (convergentes, con codo exterior) —
\draw[color=unamAzul!70] (api.south) -- ++(0,-0.55) -| (scripts.north west);
\draw[color=unamAzul!70] (vft.south) -- ++(0,-0.55) -| (scripts.north east);

%% — Flujo interno —
\draw[color=unamAzul!70] (scripts) -- (gpkg);
\draw[color=unamAzul!70] (gpkg)    -- (qgis);
\draw[color=unamAzul!70] (qgis)    -- (exports);

%% — Exports → LaTeX —
\draw[color=unamOro] (exports) -- (tesis);

%% — Recuadro del repositorio Transport-gis (en fondo) —
\begin{scope}[on background layer]
  \node[draw=unamAzul!40, dashed, rounded corners=6pt, fill=unamAzul!3,
        fit={(scripts)(gpkg)(qgis)(exports)},
        inner sep=14pt] (tgbox) {};
\end{scope}
%% — Label DENTRO del recuadro, esquina sup-izq —
\node[font=\small\bfseries\color{unamAzul}, anchor=north west]
  at ($(tgbox.north west) + (4pt,-3pt)$) {Transport-gis-zmvm-mjg};

\end{tikzpicture}
